\section{Related work}
\label{sec:related}

\paragraph{Neural tangent kernels and spectral bias.} The NTK
\cite{jacot2018neural} describes gradient descent on wide networks as kernel
regression; wide networks evolve as their linearisation \cite{lee2019wide},
and their priors are Gaussian processes \cite{neal1996bayesian, lee2018deep}.
Directions with large kernel eigenvalues are learned first --- the spectral or
frequency bias \cite{rahaman2019spectral}. Output scaling moves a network
between this lazy regime and feature learning \cite{chizat2019lazy}. We use
all of these, but on a kernel defined over \emph{pairs} and decomposed by
graph topology rather than by frequency.

\paragraph{Hodge theory for rankings.} HodgeRank
\cite{jiang2011statistical} fits the gradient component by least squares and
reads the curl and harmonic residuals as measures of inconsistency; later work
uses it for crowdsourced ranking and Bayesian intransitive models
\cite{xu2018hodgerank}. No neural network or training dynamics is involved.

\paragraph{Intransitive preference models.} Bradley--Terry
\cite{bradley1952rank} underlies RLHF reward models \cite{ouyang2022training}.
GPM \cite{zhang2025beyond} embeds responses and scores pairs with a
skew-symmetric bilinear form; HRC \cite{hrc2026transitivity} decomposes
preferences into an explicit transitive scalar and a cyclic vector part. These
works establish expressivity and downstream gains; we study the learning
dynamics that any such model is subject to.

\paragraph{NTK for pairwise feedback.} Neural dueling bandits analyse
Bradley--Terry preference learning in the NTK regime to obtain confidence
bounds \cite{verma2024neural}; they do not decompose the kernel by Hodge
structure. The closest work is the Topological NTK \cite{krishnagopal2026topological},
which reports a ``topological spectral bias'' for message passing on edge
signals of simplicial complexes; it concerns graph neural networks on a fixed
complex, not preference models over item features, and does not address
Bradley--Terry losses, annotator noise or early stopping.
